\documentclass[12pt,a4paper]{article}

% Pakieki do obsługi języka polskiego i kodowania
\usepackage[utf8]{inputenc}
\usepackage[T1]{fontenc}
\usepackage[polish]{babel}

% Dodatkowe pakiety zwiększające czytelność dokumentu
\usepackage{geometry}
\geometry{top=2.5cm, bottom=2.5cm, left=2.5cm, right=2.5cm}
\usepackage{hyperref}
\usepackage{booktabs}
\usepackage{listings}
\usepackage{xcolor}
\usepackage{graphicx} % Pakiet wymagany do wklejania obrazków (np. JPG, PNG)

% Konfiguracja wyglądu bloków kodu źródłowego
\lstset{
    language=[LaTeX]TeX,
    backgroundcolor=\color{gray!10},
    basicstyle=\small\ttfamily,
    breaklines=true,
    frame=single,
    rulecolor=\color{gray!30}
}

\title{Dokumentacja Projektu:\\System Zarządzania Kliniką Weterynaryjną}
\author{Zespół Projektowy (2 osoby)}
\date{\today}

\begin{document}
\section{Opis celu projektu}

\subsection{Problem, który rozwiązuje program}
Głównym celem projektu jest zastąpienie tradycyjnej, papierowej dokumentacji medycznej w gabinecie weterynaryjnym formą cyfrową. Program pozwala na szybkie rejestrowanie zwierząt i prowadzenie przejrzystej historii ich wizyt.

\subsection{Adresaci oprogramowania}
Program skierowany jest do:
\begin{itemize}
    \item Lekarzy weterynarii oraz personelu pomocniczego w małych i średnich lecznicach dla zwierząt,
    \item Studentów medycyny weterynaryjnej poszukujących prostego narzędzia do symulacji prowadzenia kartotek medycznych.
\end{itemize}

\subsection{Główne funkcjonalności}
Aplikacja została zaprojektowana w modelu obiektowym, oferując następujące możliwości:
\begin{enumerate}
    \item \textbf{Rejestracja pacjenta:} Dodawanie nowego zwierzęcia do bazy wraz z określeniem jego unikalnego imienia, gatunku, wieku oraz danych opiekuna.
    \item \textbf{Zarządzanie wizytami i kompozycja procedur:} Możliwość rejestracji wizyt lekarskich, gdzie pod jedną wizytę personel może przypisać nieskończenie wiele procedur medycznych.
    \item \textbf{Przegląd kartoteki (Historia wizyt):} Wyświetlanie szczegółowego wykazu wizyt oraz zabiegów dla konkretnego zwierzęcia, na podstawie wyszukania jego imienia.
    \item \textbf{Trwałość danych (Zapis/Odczyt):} Manualny zapis struktury powiązań obiektowych do pliku \texttt{dane\_kliniki.txt} przy użyciu systemu znaczników.
\end{enumerate}


% --- NOWA SEKCJA JAKO ODDZIELNY PODPUNKT ---
\newpage
\section{Diagram klas UML}

\begin{figure}[htbp]
    \centering
    \includegraphics[width=0.8\textwidth]{diagram.jpg} % Szerokość dopasowana do 80% obszaru tekstu
    \label{fig:diagram_klas_uml}
\end{figure}
\newpage
% -------------------------------------------


\section{Instrukcja obsługi (Menu główne)}

Po uruchomieniu aplikacji w oknie konsoli pojawi się interaktywne menu tekstowe. Nawigacja odbywa się poprzez wpisanie cyfry od 1 do 6 i zatwierdzenie klawiszem \texttt{Enter}:
\begin{itemize}
    \item \textbf{Opcja 1:} Pozwala zarejestrować zwierzę. Jeśli wprowadzony wiek nie będzie liczbą całkowitą dodatnią, program poprosi o poprawienie danych.
    \item \textbf{Opcja 2:} Służy do zapisu nowej wizyty. Najpierw podaj imię zarejestrowanego pacjenta. Następnie wpisz nazwisko lekarza oraz datę. Program uruchomi wewnętrzną pętlę, pytając o dane kolejnych procedur. Wpisanie litery \texttt{t} dodaje kolejną procedurę, natomiast każda inna odpowiedź kończy edycję wizyty.
    \item \textbf{Opcja 3:} Wyświetla pełną historię leczenia wybranego zwierzaka.
\item \textbf{Opcja 4:} Zapisuje aktualne dane pacjentów, ich wizyt oraz procedur do pliku tekstowego, dzięki czemu nie znikną one po wyłączeniu programu.
    \item \textbf{Opcja 5:} Wczytuje wcześniej zapisane dane z pliku tekstowego z powrotem do programu.
    \item \textbf{Opcja 6:} Bezpiecznie zamyka aplikację.
\end{itemize}


\end{document}